\section{Podsumowanie i wnioski}
\akapit Porównując wyniki uzyskane z obu algorytmów nawigacyjnych, można stwierdzić, że obie metody poprawnie odwzorowują charakter ruchu obiektu, natomiast w przypadku równań kinematycznych widać większą podatność na błędy wynikające z numerycznego całkowania i orientacji osi. Algorytm kwaternionowy charakteryzuje się płynniejszymi przejściami. Zauważalny dryft kątów w obu przypadkach jest zjawiskiem naturalnym dla systemów bezwładnościowych pozbawionych zewnętrznej korekcji (np. z magnetometru lub GPS) i wynika z kumulowania się błędów pomiarowych giroskopów w czasie.